\chapter{What a Protocol Generates}
\label{ch:generated}

\index{generated artifacts}
A LeanFM project has one program, the protocol, and many outputs. This chapter is
what comes out, why each output is generated rather than written, and how the
outputs are reviewed.

\section{Diagrams}

\index{diagram!generated}
\index{provenance}
The top-level protocol's diagram is what people read first, the way an OpenAPI
viewer is where people meet an HTTP interface. It is always LeanFM's rendering of
the current text (Figure~\ref{fig:placeorder}), never a drawing of it, so it
cannot drift from the program. Every generated view also records why it exists:
the prompt that asked for it, the question it answers, its event source, the
reducer that computed it, its unit, and its boundary. A view without those is a
picture; with them it is evidence.

Large protocols are read from the top: an included protocol is a frame that
opens, and a nested \texttt{sessions} block is a service drawn once. The same
structure could drive an animated canvas, with each frame folding open for
detail.

\section{Code each role is written on}

\index{typestate}
\index{skeleton}
From the typestates of Chapter~\ref{ch:meaning}, LeanFM writes a module with one
class per state of each role, in Python or TypeScript. A method exists only where
the protocol allows the move, fills in the cause the protocol names, logs the
event in the format conformance reads, and returns the next state. The client of
\texttt{PlaceOrder}, in TypeScript:

\begin{lstlisting}[language={},basicstyle=\ttfamily\footnotesize,breaklines=true]
export class ClientS0 extends State {
  readonly state = "ClientS0" as const;

  sendCreateOrder(body: CreateOrder): ClientS1 {
    const ep = this.take();
    ep.send("Shop", "CreateOrder", "o", null, body);
    return new ClientS1(ep);
  }
}

/** Client at S1: receive(msg): Quote from Shop, SoldOut from Shop. */
export class ClientS1 extends State {
  readonly state = "ClientS1" as const;

  receive(msg: Message): { label: "Quote"; msg: Message; body: Quote; next: ClientS2 } | { label: "SoldOut"; msg: Message; body: SoldOut; next: End } {
    const ep = this.ep;
    if (ep.accepts(msg, "Shop", "Quote")) {
      this.take();
      const body = ep.received(msg, "q") as Quote;
      return { label: "Quote", msg, body, next: new ClientS2(ep) };
    }
    if (ep.accepts(msg, "Shop", "SoldOut")) {
      this.take();
      const body = ep.received(msg, null) as SoldOut;
      return { label: "SoldOut", msg, body, next: new End(ep) };
    }
    return ep.refuse("ClientS1", msg, "Quote from Shop, SoldOut from Shop");
  }
}

/** Client at S2: sendPay -> ClientS3. */
export class ClientS2 extends State {
  readonly state = "ClientS2" as const;

  sendPay(body: Pay): ClientS3 {
    const ep = this.take();
\end{lstlisting}


In TypeScript a move the state does not allow does not type-check, and
\texttt{receive} returns a union tagged by label, so handling every answer the
protocol allows is a type obligation. The implementer writes only the decisions
(which branch, what each message carries) and the transport. Because the
generated code logs every message, conformance (Chapter~\ref{ch:execution}) can
check the running program without further instrumentation. With each message it
logs exactly the fields of its body that some \texttt{where} condition reads, and
nothing else: \texttt{PlaceOrder}'s client logs the amount it pays, never the card
number in the same body.

\section{Interfaces}

\index{OpenAPI}
\index{protobuf}
A role that serves HTTP gets its interface written out as OpenAPI 3.1: one
operation per route, the label as its operation id, and every response the
protocol allows under its status. The shop's, abridged:

\begin{lstlisting}[language={},basicstyle=\ttfamily\footnotesize,breaklines=true]
{"paths":
 {"/orders":
  {"post":
   {"summary": "CreateOrder from Client",
    "responses":
    {"409":
     {"description": "SoldOut",
      "content":
      {"application/json":
       {"schema": {"$ref": "#/components/schemas/SoldOut"}}}},
     "201":
     {"description": "Quote",
      "content":
      {"application/json":
       {"schema": {"$ref": "#/components/schemas/Quote"}}}}},
    "requestBody":
    {"content":
     {"application/json":
      {"schema": {"$ref": "#/components/schemas/CreateOrder"}}}},
    "operationId": "CreateOrder"}}},
 "openapi": "3.1.0",
 "info":
 {"version": "0.1.0",
  "title": "PlaceOrder: Shop",
  "description":
  "The requests Shop receives in protocol PlaceOrder, generated by LeanFM from the session type."},
 "components":
 {"schemas":
  {"SoldOut":
   {"type": "object",
    "title": "SoldOut",
    "required": [],
    "properties": {},
    "additionalProperties": false},
   "Quote":
   {"type": "object",
    "title": "Quote",
    "required": ["price"],
    "properties": {"price": {"type": "integer", "minimum": 0}},
    "description": "The shop's price for the order, in cents.",
\end{lstlisting}


Each body's schema is its message type, with only its declared fields. The same
types are written out as protobuf, to fix the shape and bytes of each message so
every role and every stored record agrees on them:

\begin{lstlisting}[language={},basicstyle=\ttfamily\footnotesize,breaklines=true]
// What the client orders.
message CreateOrder {
  string item = 1;
  uint64 quantity = 2;
}

// The shop's price for the order, in cents.
message Quote {
  uint64 price = 1;
}

// The client pays; the card number is a secret, never logged.
message Pay {
  uint64 amount = 1;
  string card_secr = 2;
}

message Charge {
  uint64 amount = 1;
  string card_secr = 2;
}
\end{lstlisting}


and as Lean structures that convert to and from JSON, for code written in Lean:

\begin{lstlisting}[language={},basicstyle=\ttfamily\footnotesize,breaklines=true]
/-- What an event carries: one message of the protocol, or an event a role performs alone. -/
inductive Payload where
  | createOrder (m : CreateOrder)
  | quote (m : Quote)
  | pay (m : Pay)
  | charge (m : Charge)
  | charged (m : Charged)
  | receipt (m : Receipt)
  | declined (m : Declined)
  | paymentDeclined (m : PaymentDeclined)
  | soldOut (m : SoldOut)
  deriving Repr, BEq

/-- The event envelope (selfhost.session.envelope): ts in nanoseconds, unique in its stream. -/
structure Event where
  ts : Nat
  session : String
  src : Role
  dst : Role
  cause : Option Nat := none
  prev : Option Nat := none
  /-- the sender's clock when it sent and the receiver's when it received, ns (selfhost.session.timelines) -/
  sentAt : Option Nat := none
\end{lstlisting}


The skeletons carry them too: a Python dataclass or a TypeScript interface per
message, and each send method takes its message's type. None of these say anything
about order or about who decides; that is the protocol's job. The LeanFM server's
own messages are an example (Appendix~\ref{ch:web-protocol}).

\section{Requirements and implementation, reviewed in git}

\index{Requirements.lean}
\index{Implementation.lean}
\index{review}
The durable record of a project is two Lean files written out from the checked
protocol, and the work proceeds in three layers, each checked before the next:

\begin{enumerate}
  \item \textbf{Requirements} (\texttt{Requirements.lean}): the observable
        behaviour, one clause per behaviour, each with what must always hold and
        what must never happen.
  \item \textbf{Implementation mapping} (\texttt{Implementation.lean}): how each
        clause is realized and which test or proof covers it, including every
        negative requirement.
  \item \textbf{Source}, written in one pass from the two.
\end{enumerate}

Review is what a team already does: read the diff and commit it. LeanFM records no
approvals, and a passing check is evidence for the review, not the review itself.
Nothing lives only in a chat: the protocol, the two Lean files, the generated
views and the source are all in the project's one directory, so a new session
resumes from the files. LeanFM is itself such a project; its own requirements are
\texttt{LeanFM/LLMGenerated/SelfHosting.lean}, and every feature in this book is a
clause there with an implementation decision and a test.

Chapter~\ref{ch:feature-discipline} is the discipline that keeps those clauses
honest: no feature is complete until its forbidden behaviours are stated and each
has a control and a negative test.

\section{Building software from a protocol}

\begin{enumerate}
  \item Write the protocol and check it until it has no issues.
  \item Show it: the diagram, each role's typestate, a walk through one session,
        and a generated stream.
  \item State the assumptions and promises, and check the promises against the
        assumptions.
  \item Generate each role's states and implement the role on them.
  \item Run it, and check its log for conformance and its promises over closed
        windows; fix the code, or the protocol if the protocol was wrong.
  \item State each rule the structure does not show as a CTL question.
  \item When the protocol changes, ask which roles can be redeployed first.
\end{enumerate}
